\smallskip

Dans l'espace rapporté à un repère orthonormé \Oijk, on considère les points :
\begin{center}
A$(4~;~-4~;~4)$, \quad B$(5~;~-3~;~2)$, \quad C$(6~;~- 2;~;~3)$, \quad D(5~;~1~;~1)
\end{center}

\begin{enumerate}
	\item On a : $\vect{AB} = \begin{pmatrix}5-1\\-3-(-4)\\2-4\end{pmatrix} = \begin{pmatrix}1\\1\\-2\end{pmatrix}$ et $\vect{CB} = \begin{pmatrix}5-6\\-3-(-2)\\2-3\end{pmatrix} = \begin{pmatrix}-1\\-1\\-1\end{pmatrix}$.

Donc :
$\vect{AB} \cdot \vect{CB} = 1 \times (-1) + 1 \times (-1) - 2 \times (-1) = -1 -1 + 2 = 0.$

\smallskip

Puisque leur produit scalaire est nul, les vecteurs $\vect{AB}$ et $\vect{CB}$ sont orthogonaux, donc le triangle ABC est rectangle en B.

	\item D'après la question précédente, les points A, B et C étant non alignés, ils définissent bien un plan. Leurs coordonnées vérifient donc toute équation cartésienne du plan ainsi défini.

Or on a :
$4-(-4)-8 = 4 + 4 - 8 = 0,$
donc les coordonnées de A vérifient l'équation $x - y  - 8 = 0$.
	
De même
$5 - (-3) - 8 = 5 + 3 - 8 = 0 \qquad \text{et} \qquad 6 - (-2) - 8 = 6 + 2 - 8 = 0,$
donc les coordonnées de B et C vérifient également l'équation $x - y - 8 = 0$.

Donc \quad $x - y - 8 = 0$ \quad est bien une équation cartésienne du plan (ABC).

\item
	\begin{enumerate}
		\item On sait que si $\vectt{n} = \begin{pmatrix}a\\b\\c\end{pmatrix}$ est un vecteur normal à un plan, alors ce plan a une équation cartésienne de la forme $ax + by + cz + d = 0$. Or on a vu à la question précédente que $x - y - 8 = 0$ est une équation cartésienne du plan (ABC), donc $\vectt{n} = \begin{pmatrix}1\\-1\\0\end{pmatrix}$ est un vecteur normal au plan (ABC).

Puisque $d$ est orthogonale au plan (ABC), le vecteur $\vectt{n} = \begin{pmatrix}1\\-1\\0\end{pmatrix}$ est un vecteur directeur de $d$. De plus, la droite $d$ passe par le point $\mathrm{D}\,(5~;~1~;~1)$. \\
Donc une représentation paramétrique de $d$ est :

\[\boxed{\left\{
\begin{array}{rcl}
x &=& 5 + t\\
y &=& 1 - t\\
z &=& 1
\end{array}\right. \qquad t \in \R.}\]

		\item Puisque la droite $d$ est orthogonale au plan (ABC) et passe par D, le projeté orthogonal H de D sur (ABC) est le point d'intersection de la droite $D$ et du plan (ABC). Les coordonnées $(x~;~y~;~z)$ du point H vérifient donc à la fois la représentation paramétrique de $d$ et l'équation cartésienne de (ABC). Elles sont donc solution du système :
		\[\begin{array}{l}
\left\{
\begin{array}{rcl}
x &=& 5 + t\\
y &=& 1 - t\\
z &=& 1\\
x - y - 8 &=& 0
\end{array}\right.\\
\\
\Longrightarrow 5 + t - (1 - t) - 8 = 0 \iff 2t - 4 = 0 \iff 2t = 4 \iff t = 2\\
\\
\left\{
\begin{array}{rcl}
x &=& 5 + 2 = 7\\
y &=& 1 - 2 = -1\\
z &=& 1
\end{array}\right.\end{array}\]
Les coordonnées du projeté orthogonal H de D sur (ABC) sont donc :

\[\boxed{\mathrm{H}\,(7~;~-1~;~1)}\]

		\item \quad $\mathrm{DH} = \sqrt{(7-5)^2 + (-1-1)^2 + (1 - 1)^2} = \sqrt{2^2 + (-2)^2 + 0^2} = \sqrt{4+4} = \sqrt{8} = \boxed{2\sqrt{2}}$
	\end{enumerate}

\item
	\begin{enumerate}
		\item La pyramide ABCD admet pour base le triangle ABC et pour hauteur le segment [DH].
		
		On a donc $V = \dfrac{1}{3} \times \text{aire}\,(\mathrm{ABC}) \times \mathrm{DH}$.

Puisque ABC est rectangle en B, on a :
\begin{eqnarray*}
\text{aire}\,(\text{ABC}) &=& \dfrac{\mathrm{AB}\times\mathrm{CB}}{2} \\
&=& \dfrac{\|\vect{AB}\| \times \|\vect{CB}\|}{2}\\
&=& \dfrac{\sqrt{1^2 + 1^2 + (-2)^2} \times \sqrt{(-1)^2 + (-1)^2 + (-1)^2}}{2}\\
&=& \dfrac{\sqrt{1+1+4}\times\sqrt{1+1+1}}{2} = \dfrac{\sqrt{6}\times\sqrt{3}}{2} = \dfrac{\sqrt{18}}{2}\\
&=& \dfrac{3\sqrt{2}}{2}.
\end{eqnarray*}
On a donc :

\[V = \dfrac{1}{\cancel{3}} \times \dfrac{\cancel{3}\sqrt{2}}{\cancel{2}} \times \cancel{2}\sqrt{2} = \boxed{2}\]

Le volume de la pyramide ABCD est $2$.
		\item La distance du point A au plan (BCD) est la longueur de la hauteur issue de A de la pyramide ABCD. Soit $x$ cette distance. On a alors $V = \dfrac{1}{3} \times \text{aire\,(BCD)} \times x$. Donc :

\[\dfrac{1}{3} \times \dfrac{\sqrt{42}}{2} \times x = 2 \iff x\sqrt{42} = 2 \times 3 \times 2 \iff x = \dfrac{12}{\sqrt{42}} = \dfrac{2\times6}{\sqrt{6} \times \sqrt{7}} = \dfrac{2\sqrt{6}}{\sqrt{7}} = \boxed{\dfrac{2\sqrt{42}}{7}}\]
La distance du point A au plan (BCD) est $\dfrac{2\sqrt{42}}{7}$.
	\end{enumerate}
\end{enumerate}


\bigskip
